\section{Results}
\label{sec:results}

\subsection{Model Availability}
Of the nineteen models registered in the system, twelve were successfully trained and evaluated on the 366-observation evaluation dataset described in Section~\ref{sec:dataset}: all seven statistical models, all four classical machine learning models, and the Croston--SBA intermittent-demand model. The seven deep learning models (RNN, LSTM, GRU, BiLSTM, CNN-1D, TCN, and Transformer) were withheld automatically by the frequency-aware gating mechanism described in Section~\ref{subsec:dl-models}, since the dataset's 366 daily observations fall below the system's configured minimum of 500 observations for reliable deep learning training at daily frequency; this decision, and its stated justification, was returned directly in the system's response rather than only logged internally.

\subsection{Quantitative Comparison}
Table~\ref{tab:model-comparison} reports the full leaderboard of all twelve successfully trained models, ranked by the composite score $S$ of Eq.~\eqref{eq:composite-score}, together with each model's cross-validated RMSE, MAE, MAPE, $R^2$, and mean per-fold training and prediction time. Figure~\ref{fig:rmse-comparison} presents the same RMSE values sorted independently by magnitude for direct visual comparison. The $k$-nearest-neighbor regressor was selected as the best model by the composite score, with a score of 0.7983, despite not attaining the lowest RMSE among the twelve candidates; XGBoost attained the lowest raw RMSE (1109.45) but ranked ninth overall (score 0.9460), a divergence discussed in detail in Section~\ref{sec:discussion}.

\begin{table*}[!t]
\caption{Full Model Leaderboard on the Evaluation Dataset (366 Daily Observations), Ranked by Composite Score $S$ (Eq.~\eqref{eq:composite-score})}
\label{tab:model-comparison}
\centering
\begin{tabular}{@{}clrrrrrrr@{}}
\toprule
\textbf{Rank} & \textbf{Model} & \textbf{Score $S$} & \textbf{RMSE} & \textbf{MAE} & \textbf{MAPE} & \textbf{$R^2$} & \textbf{Train (s)} & \textbf{Predict (s)} \\
\midrule
1  & KNN                  & 0.7983 & 1162.41 & 907.53 & 3.002 & $-0.078$ & 0.0001 & 0.2010 \\
2  & SVM                  & 0.8036 & 1125.86 & 894.23 & 3.213 & $\phantom{-}0.003$ & 0.0016 & 0.0006 \\
3  & Exponential Smoothing & 0.8221 & 1129.37 & 882.53 & 3.368 & $-0.011$ & 0.0014 & 0.0006 \\
4  & Moving Average       & 0.8261 & 1158.52 & 897.45 & 3.400 & $-0.062$ & 0.0000 & 0.0004 \\
5  & Croston--SBA         & 0.8271 & 1141.70 & 891.96 & 3.349 & $-0.036$ & 0.0002 & 0.0000 \\
6  & Holt--Winters        & 0.8387 & 1144.76 & 902.69 & 3.397 & $-0.040$ & 0.1764 & 0.0043 \\
7  & Prophet              & 0.8608 & 1215.07 & 923.64 & 2.944 & $-0.168$ & 0.3744 & 0.0080 \\
8  & SARIMA               & 0.8780 & 1149.83 & 900.94 & 3.377 & $-0.048$ & 0.1107 & 1.1075 \\
9  & XGBoost              & 0.9460 & 1109.45 & 865.96 & 3.246 & $\phantom{-}0.026$ & 1.9747 & 0.0081 \\
10 & SARIMAX              & 0.9966 & 1325.73 & 1020.74 & 3.251 & $-0.413$ & 0.2205 & 1.0920 \\
11 & ARIMA                & 1.0161 & 1130.18 & 894.03 & 3.517 & $-0.014$ & 0.1052 & 4.1640 \\
12 & Random Forest        & 1.0810 & 1224.23 & 965.23 & 3.995 & $-0.221$ & 1.8084 & 0.3019 \\
\bottomrule
\end{tabular}
\end{table*}

% RMSE comparison bar chart (real values from Table IV / model_comparison.tex)
\begin{figure}[!t]
\centering
\begin{tikzpicture}[scale=1]
\foreach \name/\val [count=\i from 0] in {
    XGBoost/1109.45,
    SVM/1125.86,
    Ex.Smooth/1129.37,
    ARIMA/1130.18,
    CrostonSBA/1141.70,
    Holt-Winters/1144.76,
    SARIMA/1149.83,
    Moving Avg/1158.52,
    KNN/1162.41,
    Prophet/1215.07,
    RandomForest/1224.23,
    SARIMAX/1325.73
} {
    \pgfmathsetmacro{\ypos}{-0.4*\i}
    \pgfmathsetmacro{\barlen}{(\val-1050)/40}
    \node[anchor=east, font=\tiny] at (-0.15,\ypos) {\name};
    \draw[fill=orange!35, draw=orange!60] (0,\ypos-0.13) rectangle (\barlen,\ypos+0.13);
    \node[anchor=west, font=\tiny] at (\barlen+0.05,\ypos) {\pgfmathprintnumber[fixed,precision=0]{\val}};
}
\draw[-{Latex[length=2mm]}] (0,-5.1) -- (7.2,-5.1) node[right, font=\tiny]{RMSE};
\end{tikzpicture}
\caption{Cross-validated RMSE (Eq.~\eqref{eq:rmse}) across all twelve successfully trained models on the evaluation dataset, sorted ascending. Note that XGBoost attains the lowest raw RMSE but ranks ninth under the composite score $S$ of Eq.~\eqref{eq:composite-score} (Table~\ref{tab:model-comparison}), discussed in Section~\ref{sec:discussion}.}
\label{fig:rmse-comparison}
\end{figure}


\subsection{Coefficient of Determination}
Ten of the twelve successfully trained models produced a negative cross-validated $R^2$ on this dataset, and the two positive values (support vector regression at 0.003 and XGBoost at 0.026) are close to zero. A negative $R^2$ indicates that the model's cross-validated forecast performed worse, in squared-error terms, than a naive forecast equal to the mean of the training window for that fold; this is a genuine and honestly reported characteristic of this specific dataset's difficulty, discussed further in Section~\ref{sec:discussion}, rather than an error in the evaluation methodology.

\subsection{Mean Absolute Percentage Error}
The reported MAPE values (Eq.~\eqref{eq:mape}) range from approximately 2.94 to 3.99, that is, 294\% to 399\%, substantially exceeding the conventional expectation that MAPE be expressed as a small fraction. This arises from the definition of MAPE itself, which is undefined at, and highly sensitive near, true values close to zero; the evaluation dataset's daily aggregated sales include days with very low transaction volume (Table~\ref{tab:dataset-stats} reports a minimum daily value of 25.00 against a mean of 1{,}299.89), and a small absolute error on such a low-value day contributes a disproportionately large percentage error. This is reported transparently as an observed limitation of the MAPE metric for this dataset rather than adjusted after the fact.

\subsection{Feature Importance}
For the gradient boosting model, Figure~\ref{fig:feature-importance} reports the top eight gain-based feature importances obtained directly from the trained model on this dataset. The most influential features are a 14-day rolling maximum, a 7-day exponentially weighted moving average, and calendar-cyclical encodings of month, indicating that the model relies primarily on recent local history and monthly seasonality rather than any single dominant signal, consistent with the absence of exogenous price or promotional columns in this particular evaluation dataset.

% Real feature-importance bar chart (XGBoost, evaluation dataset, Section VI)
\begin{figure}[!t]
\centering
\begin{tikzpicture}[scale=1]
\foreach \name/\val [count=\i from 0] in {
    roll\_max\_14/0.0964,
    ewm\_7/0.0711,
    cos\_month/0.0682,
    lag\_21/0.0605,
    month/0.0600,
    lag\_1/0.0543,
    roll\_max\_7/0.0514,
    lag\_28/0.0505
} {
    \pgfmathsetmacro{\ypos}{-0.5*\i}
    \pgfmathsetmacro{\barlen}{\val*40}
    \node[anchor=east, font=\scriptsize] at (-0.15,\ypos) {\texttt{\name}};
    \draw[fill=blue!35, draw=blue!60] (0,\ypos-0.15) rectangle (\barlen/10,\ypos+0.15);
    \node[anchor=west, font=\tiny] at (\barlen/10+0.05,\ypos) {\val};
}
\draw[-{Latex[length=2mm]}] (0,-4.1) -- (4.2,-4.1) node[right, font=\tiny]{gain-based importance};
\end{tikzpicture}
\caption{Top eight gain-based feature importances from the XGBoost model fit on the evaluation dataset of Section~\ref{sec:dataset} (real values reported by the trained model, not illustrative data). Calendar features (\texttt{cos\_month}, \texttt{month}) and short/medium-range history (\texttt{lag\_1}, \texttt{lag\_21}, \texttt{lag\_28}, rolling maxima) dominate, consistent with the absence of strong exogenous signal in this particular dataset.}
\label{fig:feature-importance}
\end{figure}


\subsection{Uncertainty Intervals}
The selected model (KNN) received a $95\%$ forecast interval derived from its own historical cross-validation RMSE under a normal approximation widening with the square root of the forecast step, following the general procedure described in Section~\ref{sec:proposed}; the random forest model, on the seven occasions it is the best model on other datasets exercised during system development, instead receives an interval derived from the empirical spread of its individual trees' predictions, since this is a strictly more informative interval than the generic RMSE-based fallback when it is available.

\subsection{Computation Time}
Training and prediction times in Table~\ref{tab:model-comparison} are per-fold averages under the cross-validation procedure of Section~\ref{subsec:cv}, measured on the hardware described in Section~\ref{sec:experimental}. Instance-based and simple statistical methods (KNN, SVM, moving average, exponential smoothing, Croston--SBA) trained in under two milliseconds per fold, reflecting their lack of an iterative optimization procedure, while ARIMA's grid-searched order selection and the tree ensembles' iterative fitting incurred training times between roughly 0.1 and 2 seconds per fold; ARIMA's per-fold \emph{prediction} time (4.164~s) was the largest recorded value in this experiment, arising from one-step-ahead walk-forward re-estimation across the held-out fold, described in Section~\ref{subsec:cv}. All twelve models nonetheless completed within their configured per-category timeout, and the complete nineteen-model submission (including the seven gated-out deep learning models, which return immediately without training) completed end-to-end, from receipt of the forecast request to a fully populated leaderboard response, in 18.6 seconds for this dataset and a 30-day forecast horizon.
